\section{Introduction}
\label{sec:introduction}

Classical centrally planned economies relied on static material balance
techniques that struggled to incorporate dynamic consumer preferences,
capital accumulation, and structural change
\citep{montias1959planning,ellman1989socialist}.
As \citet{kornai1992socialist} documented, these systems suffered from
chronic shortage, soft budget constraints, and inability to respond to
changing demand patterns.
Conversely, decentralized market systems utilize price signals to
coordinate production but suffer from persistent instability,
inequality, and inefficient investment dynamics
\citep{kalecki1971selected,minsky1986stabilizing}.

The socialist calculation debate of the 1920s--1940s centred on whether
rational economic calculation was possible without market prices
\citep{mises1920economic,hayek1945use,lange1936economic}.
\citet{lange1936economic} and \citet{taylor1929guidance} proposed a
``trial-and-error'' procedure using shadow prices, anticipating modern
computational approaches.
However, their proposals remained theoretical due to computational
limitations.
\citet{kantorovich1965best} demonstrated the feasibility of linear
programming for optimal planning in the U.S.S.R. introducing
shadow prices derived from dual variables.
\citet{leontief1986input} developed dynamic input-output models that
could track inter-industry flows over time, though growth remained
largely exogenous in his framework.

Recent advances in computational economics and control theory
\citep{doyle1992feedback,astrom2010feedback} make possible a new
generation of cybernetic planning models that integrate optimization,
dynamic input-output analysis, and adaptive feedback.
\citet{cottrell1993towards} proposed using linear programming and
computers for economic planning, while \citet{cockshott2012economic}
developed detailed computational algorithms for resource allocation.
\citet{allin2020economic} extended these ideas with modern machine
learning techniques.

This paper contributes to this literature by presenting a
Lagrangian-based cybernetic planning system with the following novel
features:
%
\begin{enumerate}
    \item \textbf{Endogenous growth inference}: Growth rates are derived
          from revealed consumption demand following
          \citet{samuelson1948consumption}, not imposed administratively.
          
    \item \textbf{Shadow price feedback}: Prices derived from constrained
          utility maximization serve as informational signals for
          decentralized adjustment, addressing
          \citeauthor{hayek1945use}'s (\citeyear{hayek1945use}) knowledge
          problem.

    \item \textbf{Adaptive preferences}: Preference weights evolve based
          on observed expenditure patterns, enabling demand-led structural
          change.

    \item \textbf{Two-loop architecture}: Fast monthly price adjustments
          and slow quarterly growth planning operate simultaneously,
          inspired by hierarchical control systems
          \citep{mesarovic1970theory}.

    \item \textbf{HANK Benchmarking}: The system is rigorously validated against 
          a Heterogeneous Agent New Keynesian (HANK) benchmarking model. 
          Monte Carlo results demonstrate that while the HANK model produces 
          higher raw GDP growth (median 16.2\% vs 12.6\%), the cybernetic planner 
          yields a 10.9\% welfare advantage by resolving the allocative 
          inefficiencies and unconstrained investment cycles inherent in 
          decentralized market coordination.
\end{enumerate}
